% SOURCE_AUTHORITY: sga5_fr_workpass.tex, lines 12683--12695 (LF-normalized unit).
% SOURCE_SHA256: 791F4EFFC5E02832D5D77ED03518C8156D6F07E4C8238B03545DB93D883FBB28
\subsubsection*{4.6}

Os invariantes $Sw_y$, $Sw'_y$, $Art_y$ podem ser definidos em condições mais gerais. Para isso, seja $C'$ uma curva algébrica lisa sobre $k$, não necessariamente conexa. Suponha-se que o grupo finito $G$ opere sobre $C'$ e que exista um aberto $U$ de $C=C'/G$, denso em $C$, tal que $C'|U$ seja um recobrimento principal de $U$ de grupo $G$. Para um ponto $y\in C$, podem-se definir os módulos $Sw_y$, $Sw'_y$, $Art_y$ procedendo como antes. Pode-se reduzir ao caso em que $C'$ é conexa da seguinte maneira. Sejam $y'\in C'$ com $p(y')=y$, $C'_1$ a componente irredutível de $C'$ que contém $y'$, e $C_1$ a componente irredutível de $C$ que contém $y$. Se $G_{C'_1}$ é o estabilizador da componente $C'_1$ e se $\eta'_1$ (resp. $\eta_1$) é o ponto genérico de $C'_1$ (resp. $C_1$), então a extensão
\[
K'_1=k(\eta'_1)\supset k(\eta_1)=K_1
\]
é de Galois, com grupo de Galois $G_{C'_1}$, e $C'_1$ é a normalização de $C_1$ em $K'_1$ (SGA 1 I 10). Podem-se aplicar, portanto, as considerações anteriores às curvas $C'_1$, $C_1$ e ao grupo $G_{C'_1}$. Desse modo, obtêm-se os $\mathbb Z_\ell[G_{C'_1}]$-módulos projetivos de tipo finito
\[
Sw_{G_{C'_1},y},\qquad Sw'_{G_{C'_1},y},\qquad Art_{G_{C'_1},y},
\]
e $Sw_y$, $Sw'_y$, $Art_y$ são os módulos induzidos associados à inclusão $G_{C'_1}\subset G$.

\textit{Observação 4.6.} Os módulos $Sw_{y'}$, $Sw'_{y'}$, $Art_{y'}$ são invariantes puramente locais: são determinados pelo grupo $G_{y'}$ que opera sobre o completamento do anel local $\mathcal O_{C',y'}$. Além disso, $Sw_{y'}=0$ se, e somente se, $p:C'\to C$ é moderadamente ramificado em $y'$, ou, o que é equivalente, se a ordem de $G_{y'}$ é prima com a característica de $k$ [3 VI].
